\section{Riesgos técnicos}
Los riesgos técnicos de Kit Pagos Colombia surgen principalmente de las consecuencias negativas ya identificadas en los ADR de este documento, como el desfase entre los mocks del simulador y las respuestas reales de las pasarelas, la cobertura desigual de la verificación de firmas entre los cuatro manejadores de webhook, o la clasificación de errores transitorios en el mecanismo de reintentos, junto con riesgos propios de la infraestructura de despliegue y del alcance limitado por el tiempo del proyecto académico.

En resumen, los riesgos de prioridad alta se concentran en la fidelidad de los mocks, el arranque en frío del plan gratuito de Render, la clasificación de errores del reintento, el mantenimiento de la tabla de estados y la validación de las firmas de webhook en producción. Los riesgos de prioridad media y baja cubren la actualización manual de la fábrica de adaptadores, la latencia acumulada por reintentos y el cambio de condiciones comerciales de un proveedor externo, este último ya materializado con la transición de PayU a Rapyd.

Dado que cada riesgo requiere documentar probabilidad, impacto, prioridad, plan de mitigación y plan de contingencia, presentar esta información en una tabla dentro del documento obligaría a comprimir el texto o reducir el tamaño de la fuente para que fuera legible. Por esta razón, al igual que con los ASR, se optó por adjuntar una hoja de cálculo con el detalle completo de los riesgos técnicos: \href{https://docs.google.com/spreadsheets/d/1eXwtLcUxgRN3zpxN8TdkdVaR1zTJFfTZaZW77_ftcfA/edit?usp=sharing}{Riesgos técnicos Kit Pagos Colombia}.

